% =====================================================================
%  Manual técnico — ServiceDesk UTA
%  Compilar:  latexmk -pdf manual_tecnico.tex
%  (o dos veces: pdflatex manual_tecnico.tex)
%
%  Este documento NO contiene contraseñas, secretos, direcciones de
%  servidores reales ni datos personales. Los valores entre <...> son
%  marcadores que cada instalación sustituye.
% =====================================================================
\documentclass[11pt,a4paper,oneside]{report}

% ---------- Idioma y codificación ----------
\usepackage[utf8]{inputenc}
\usepackage[T1]{fontenc}
\usepackage[spanish,es-noshorthands,es-tabla]{babel}
\usepackage{lmodern}
\usepackage{microtype}

% ---------- Página ----------
\usepackage[a4paper,top=2.5cm,bottom=2.5cm,left=2.6cm,right=2.4cm,headheight=14pt]{geometry}
\usepackage{fancyhdr}
\usepackage{titlesec}
\usepackage{enumitem}
\setlist{itemsep=2pt,topsep=4pt}

% ---------- Tablas, color, código, figuras ----------
\usepackage[table]{xcolor}
\usepackage{booktabs}
\usepackage{tabularx}
\usepackage{longtable}
\usepackage{array}
\usepackage{listings}
\usepackage[most]{tcolorbox}
\usepackage{tikz}
\usetikzlibrary{positioning,arrows.meta,fit,backgrounds}
\usepackage{graphicx}
\usepackage[hidelinks]{hyperref}
\usepackage{bookmark}

% ---------- Colores ----------
\definecolor{utaRojo}{HTML}{8B1A1A}
\definecolor{gris}{HTML}{5A5A5A}
\definecolor{grisClaro}{HTML}{F3F3F3}
\definecolor{grisLinea}{HTML}{CFCFCF}
\definecolor{azulNota}{HTML}{1F4E79}
\definecolor{ambar}{HTML}{8A5A00}

\hypersetup{colorlinks=true,linkcolor=utaRojo,urlcolor=azulNota,citecolor=utaRojo,
  pdftitle={Manual técnico — ServiceDesk UTA},pdfauthor={Equipo ServiceDesk UTA}}

% ---------- Títulos ----------
\titleformat{\chapter}[hang]{\normalfont\huge\bfseries\color{utaRojo}}{\thechapter}{0.6em}{}
\titlespacing*{\chapter}{0pt}{-10pt}{18pt}
\titleformat{\section}{\normalfont\Large\bfseries\color{utaRojo}}{\thesection}{0.6em}{}
\titleformat{\subsection}{\normalfont\large\bfseries}{\thesubsection}{0.6em}{}

% ---------- Encabezados ----------
\pagestyle{fancy}
\fancyhf{}
\fancyhead[L]{\small\color{gris}ServiceDesk UTA · Manual técnico}
\fancyhead[R]{\small\color{gris}\nouppercase{\leftmark}}
\fancyfoot[C]{\small\thepage}
\renewcommand{\headrulewidth}{0.4pt}
\fancypagestyle{plain}{\fancyhf{}\fancyfoot[C]{\small\thepage}\renewcommand{\headrulewidth}{0pt}}

% ---------- Código ----------
\lstdefinestyle{shell}{
  basicstyle=\ttfamily\small,
  backgroundcolor=\color{grisClaro},
  frame=single, rulecolor=\color{grisLinea},
  framesep=4pt, xleftmargin=6pt, xrightmargin=6pt,
  breaklines=true, breakatwhitespace=false,
  postbreak=\mbox{\textcolor{gris}{$\hookrightarrow$}\space},
  columns=fullflexible, keepspaces=true, showstringspaces=false,
  commentstyle=\color{gris}\itshape,
  literate={á}{{\'a}}1 {é}{{\'e}}1 {í}{{\'i}}1 {ó}{{\'o}}1 {ú}{{\'u}}1
           {Á}{{\'A}}1 {É}{{\'E}}1 {Í}{{\'I}}1 {Ó}{{\'O}}1 {Ú}{{\'U}}1
           {ñ}{{\~n}}1 {Ñ}{{\~N}}1 {·}{{\textperiodcentered}}1
           {→}{{$\rightarrow$}}1 {…}{{\ldots}}1
}
\lstset{style=shell}
\lstnewenvironment{bash}{\lstset{style=shell,language=bash,morecomment=[l]\#}}{}

% \cod permite cortar la línea tras cada guion bajo; \ruta, tras cada / . _ y -.
\newcommand{\codus}{\textunderscore\allowbreak}
\DeclareRobustCommand{\cod}[1]{{\small\ttfamily\let\_\codus #1}}
\makeatletter
\g@addto@macro{\UrlBreaks}{\do\_\do\-\do\.}
\makeatother
\urlstyle{tt}
\DeclareRobustCommand{\ruta}[1]{{\small\nolinkurl{#1}}}
\setlength{\emergencystretch}{2.5em}
% En los marcadores del PDF, \cod y \ruta se muestran como texto plano.
\pdfstringdefDisableCommands{\def\cod#1{#1}\def\ruta#1{#1}\def\_{\textunderscore}}

% ---------- Cuadros ----------
\newtcolorbox{nota}[1][Nota]{colback=azulNota!5,colframe=azulNota,fonttitle=\bfseries,
  title=#1,boxrule=0.6pt,arc=2pt,left=6pt,right=6pt,top=4pt,bottom=4pt}
\newtcolorbox{aviso}[1][Importante]{colback=ambar!7,colframe=ambar,fonttitle=\bfseries,
  title=#1,boxrule=0.6pt,arc=2pt,left=6pt,right=6pt,top=4pt,bottom=4pt}

% ---------- Tablas ----------
\renewcommand{\arraystretch}{1.25}
\newcolumntype{L}{>{\raggedright\arraybackslash}X}
\newcolumntype{P}[1]{>{\raggedright\arraybackslash}p{#1}}
\newcommand{\cabecera}{\rowcolor{utaRojo!12}}

\setcounter{tocdepth}{1}
\setcounter{secnumdepth}{2}

% =====================================================================
\begin{document}

% ---------- Portada ----------
\begin{titlepage}
  \centering
  \vspace*{2.5cm}
  {\color{utaRojo}\rule{\linewidth}{1.2pt}}\\[14pt]
  {\Huge\bfseries ServiceDesk UTA\par}
  \vspace{8pt}
  {\LARGE Manual técnico\par}
  \vspace{6pt}
  {\color{utaRojo}\rule{\linewidth}{1.2pt}}\\[30pt]
  {\large Sistema de mesa de servicio basado en microservicios\par}
  \vspace{4pt}
  {\large Universidad Técnica de Ambato\par}
  \vfill
  \begin{tabular}{rl}
    \textbf{Versión del sistema:} & 1.0.0 (candidata a release) \\
    \textbf{Fecha del documento:} & 9 de noviembre de 2026 \\
    \textbf{Dirigido a:} & desarrolladores y operadores del servidor \\
  \end{tabular}
  \vspace{2cm}

  {\small\color{gris}Este manual no contiene contraseñas, secretos, direcciones de servidores\\
  ni datos personales. Los valores entre \cod{<...>} son marcadores que cada instalación sustituye.\par}
\end{titlepage}

\pagenumbering{roman}
\tableofcontents
\clearpage
\pagenumbering{arabic}

% =====================================================================
\chapter{Introducción}

\section{Propósito del sistema}
ServiceDesk UTA es una mesa de servicio web y móvil para registrar, atender y cerrar
solicitudes de soporte técnico de la comunidad universitaria. Una persona crea un
\emph{ticket} con el servicio afectado, la descripción, la ubicación y fotografías.
El personal técnico lo acepta y lo resuelve. Finalmente, quien lo creó acepta o rechaza la solución
mediante una \emph{validación digital} que queda registrada de forma inmutable.

\section{Alcance de este manual}
El manual describe la arquitectura, el modelo de datos, la API, el frontend, la seguridad,
la configuración, la instalación en un servidor nuevo, la operación diaria, las copias de
seguridad y la solución de problemas. Se basa en el código y la documentación del
repositorio en la versión 1.0.0.

\begin{itemize}
  \item \textbf{Desarrolladores}: capítulos \ref{cap:arquitectura} a \ref{cap:seguridad} y \ref{cap:pruebas}.
  \item \textbf{Operadores del servidor}: capítulos \ref{cap:configuracion} a \ref{cap:problemas}.
\end{itemize}

\section{Documentación complementaria}
El repositorio contiene guías detalladas que este manual resume:

\begin{center}
\begin{tabularx}{\linewidth}{@{}P{6.3cm}L@{}}
\toprule
\cabecera \textbf{Carpeta / archivo} & \textbf{Contenido} \\
\midrule
\ruta{docs/architecture/} & Arquitectura de la aplicación y del servidor \\
\ruta{docs/installation/} & Requisitos, instalación, red, nginx, SQL Server, producción, seguridad \\
\ruta{docs/operations/} & Backups, Backblaze B2, retención, restauración, usuarios, monitoreo \\
\ruta{docs/troubleshooting/} & Problemas frecuentes \\
\ruta{docs/security/} & Estado de seguridad antes del servidor nuevo \\
\ruta{scripts/README.md} & Descripción de los scripts \\
\ruta{README.md} & Puesta en marcha rápida del laboratorio y reglas del repositorio \\
\bottomrule
\end{tabularx}
\end{center}

% =====================================================================
\chapter{Arquitectura}\label{cap:arquitectura}

\section{Visión general}
El sistema sigue una arquitectura de microservicios. Cada servicio es dueño de su propia
base de datos y no lee las tablas de otro servicio. Toda petición del navegador o de la
aplicación Android entra por un único punto, el \emph{API Gateway}.

\begin{figure}[h]
\centering
\begin{tikzpicture}[
  font=\small,
  caja/.style={draw=utaRojo,fill=white,rounded corners=2pt,minimum width=2.7cm,minimum height=0.9cm,align=center,line width=0.6pt},
  bd/.style={draw=azulNota,fill=azulNota!6,rounded corners=2pt,minimum width=2.3cm,minimum height=0.75cm,align=center},
  flecha/.style={-{Stealth[length=2.2mm]},line width=0.6pt,color=gris},
  node distance=0.55cm and 0.55cm]
  \node[caja] (cli) {Navegador /\\App Android};
  \node[caja,below=of cli] (nginx) {nginx del host\\(TLS 443)};
  \node[caja,below left=0.7cm and 0.4cm of nginx] (front) {frontend\\(nginx + React)};
  \node[caja,below right=0.7cm and 0.4cm of nginx] (gw) {api-gateway\\(Express)};
  \node[caja,below=1.1cm of gw] (cat) {catalog-service};
  \node[caja,left=0.5cm of cat] (auth) {auth-service};
  \node[caja,right=0.5cm of cat] (tic) {ticket-service};
  \node[bd,below=0.9cm of auth] (adb) {auth\_db};
  \node[bd,below=0.9cm of cat] (cdb) {catalog\_db};
  \node[bd,below=0.9cm of tic] (tdb) {ticket\_db};
  \begin{scope}[on background layer]
    \node[draw=azulNota,dashed,rounded corners=3pt,fit=(adb)(cdb)(tdb),inner sep=8pt,
      label={[font=\footnotesize\color{azulNota}]below:SQL Server 2022 (contenedor mssql)}] {};
  \end{scope}
  \draw[flecha] (cli) -- (nginx);
  \draw[flecha] (nginx) -- node[left,font=\footnotesize]{/} (front);
  \draw[flecha] (nginx) -- node[right,font=\footnotesize]{/api} (gw);
  \draw[flecha] (gw) -- (auth);
  \draw[flecha] (gw) -- (cat);
  \draw[flecha] (gw) -- (tic);
  \draw[flecha] (auth) -- (adb);
  \draw[flecha] (cat) -- (cdb);
  \draw[flecha] (tic) -- (tdb);
\end{tikzpicture}
\caption{Componentes y recorrido de una petición. Las llamadas internas entre servicios se describen en la sección~\ref{sec:interna}.}
\end{figure}

\section{Componentes}
\begin{center}
\begin{tabularx}{\linewidth}{@{}P{3.1cm}L@{}}
\toprule
\cabecera \textbf{Componente} & \textbf{Responsabilidad} \\
\midrule
\cod{frontend} & Interfaz React servida por nginx. Es también la base de la aplicación Android (Capacitor). \\
\cod{api-gateway} & Punto de entrada único. Aplica CORS y cabeceras de seguridad y reenvía cada prefijo \cod{/api/...} al servicio que corresponde. \\
\cod{auth-service} & Usuarios, roles, bloques, inicio de sesión local y con Microsoft, sesiones JWT y registro de accesos. \\
\cod{catalog-service} & Catálogo de servicios, reglas de enrutamiento, ubicaciones y base de conocimiento. \\
\cod{ticket-service} & Tickets, historial de eventos, adjuntos, asignación, escalamiento, validación digital, reportes, notificaciones por correo y push. \\
\cod{mssql} & SQL Server 2022 con las tres bases de datos. \\
\cod{mssql-init} & Contenedor de un solo uso que crea las bases y un \emph{login} propio para cada servicio. \\
\bottomrule
\end{tabularx}
\end{center}

\section{Enrutamiento del gateway}\label{sec:gateway}
\begin{center}
\begin{tabularx}{\linewidth}{@{}P{4.2cm}P{3.2cm}L@{}}
\toprule
\cabecera \textbf{Prefijo público} & \textbf{Servicio} & \textbf{Ruta interna} \\
\midrule
\cod{/api/auth/*} & auth-service & \cod{/auth/*} \\
\cod{/api/catalog/*} & catalog-service & se elimina el prefijo (\cod{/services}, \cod{/locations}) \\
\cod{/api/knowledge/*} & catalog-service & \cod{/knowledge/*} \\
\cod{/api/tickets/*} & ticket-service & \cod{/tickets/*} \\
\cod{/api/admin/*} & ticket-service & \cod{/admin/*} \\
\cod{/api/notifications/*} & ticket-service & \cod{/notifications/*} \\
\cod{/api/routing/*} & ticket-service & obsoleto (responde 410) \\
\cod{/auth/*} & auth-service & \cod{/auth/*}. Solo para el \emph{Redirect URI} de Microsoft (\cod{/auth/callback}) \\
\bottomrule
\end{tabularx}
\end{center}

Cada servicio expone \cod{/health} (proceso vivo) y \cod{/health/ready} (dependencias
listas, incluida la base de datos; en ticket-service también el volumen de imágenes). En el
gateway, \cod{/health/ready} agrega el estado de los tres servicios. Si un servicio no responde,
el gateway devuelve \textbf{502} con el mensaje «Servicio no disponible».

\section{Comunicación interna entre servicios}\label{sec:interna}
Además del tráfico que llega por el gateway, los servicios se llaman entre sí por HTTP dentro
de la red \cod{app}:
\begin{itemize}
  \item \textbf{catalog-service y ticket-service $\rightarrow$ auth-service}: cada petición autenticada se
        valida con \cod{POST /auth/verify}. Así, una sesión revocada deja de funcionar en todos los
        servicios, no solo en auth-service.
  \item \textbf{ticket-service $\rightarrow$ catalog-service}: consulta el servicio elegido, sus reglas de
        enrutamiento (área y prioridad) y los artículos de conocimiento vinculados a una solución.
  \item \textbf{ticket-service $\rightarrow$ auth-service}: obtiene la lista de técnicos y sus bloques para
        la asignación y los reportes.
\end{itemize}

\section{Arquitectura del servidor}
\begin{itemize}
  \item \textbf{Host}: Ubuntu Server 24.04 LTS (x86\_64). El cortafuegos UFW solo admite los puertos 22, 80 y 443.
  \item \textbf{nginx del host}: termina TLS con certificados de Let's Encrypt. Reenvía \cod{/} al frontend y \cod{/api/} al gateway, que escuchan únicamente en \cod{127.0.0.1}. Además:
    \begin{itemize}
      \item \cod{/auth/callback} va al gateway (retorno de Microsoft) y \cod{/healthz} expone el \cod{/health/ready} del gateway.
      \item \cod{/api/auth/login} tiene un límite de 10 peticiones por minuto por IP (ráfaga de 5); al superarlo responde 429.
      \item \cod{/uploads/} responde 404: las imágenes nunca se sirven como archivos estáticos.
      \item El tamaño máximo de una petición es de 30\,MB y se añaden cabeceras de seguridad (HSTS, \cod{Permissions-Policy}, etc.).
    \end{itemize}
  \item \textbf{Redes de Docker}:
    \begin{itemize}
      \item \cod{edge}: frontend y gateway.
      \item \cod{app}: gateway y servicios.
      \item \cod{data}: red \emph{interna} sin salida a Internet, con los servicios y SQL Server.
    \end{itemize}
    SQL Server no publica ningún puerto.
  \item \textbf{Volúmenes}:
    \begin{itemize}
      \item \cod{mssql-system}, \cod{mssql-data}, \cod{mssql-tlog} y \cod{mssql-backup}: datos y copias de SQL Server.
      \item \cod{ticket-uploads}: fotografías adjuntas.
    \end{itemize}
  \item \textbf{Orden de arranque}: \cod{mssql-volumes} $\rightarrow$ \cod{mssql} $\rightarrow$ \cod{mssql-init} $\rightarrow$ auth y catalog $\rightarrow$ ticket $\rightarrow$ gateway $\rightarrow$ frontend. Cada paso espera a que el anterior esté \emph{healthy}.
  \item \textbf{Límites de recursos}:
    \begin{itemize}
      \item SQL Server: 3\,GiB y 2 CPU.
      \item Cada microservicio: 512\,MiB y 1 CPU.
      \item Gateway: 256\,MiB y 0,5 CPU.
      \item Frontend: 128\,MiB y 0,5 CPU.
    \end{itemize}
  \item \textbf{Configuración común de los contenedores}: \cod{NODE\_ENV=production}, \cod{no-new-privileges},
        reinicio \cod{unless-stopped} y logs \cod{json-file} rotados (5 archivos de 20\,MB).
\end{itemize}

% =====================================================================
\chapter{Tecnologías y versiones}

\begin{center}
\begin{tabularx}{\linewidth}{@{}P{4.4cm}P{3.6cm}L@{}}
\toprule
\cabecera \textbf{Tecnología} & \textbf{Versión} & \textbf{Uso} \\
\midrule
Node.js & 22.23.4 LTS & Ejecución. Los 3 microservicios usan \emph{slim} (Prisma falla en alpine); gateway y compilación del frontend usan \emph{alpine} \\
TypeScript & 5.7 & Lenguaje de backend y frontend \\
Express & 4.21 & Servidores HTTP \\
http-proxy-middleware & 3.0 & Reenvío en el gateway \\
Prisma ORM & 5.22.0 & Acceso a datos y sincronización del esquema \\
SQL Server & 2022 CU27-GDR1 (fijado por digest) & Base de datos \\
React & 18.3.1 & Interfaz \\
Vite & 5.4.11 & Compilación del frontend \\
Tailwind CSS & 3.4.16 & Estilos \\
Capacitor & 8.5.1 (push 8.1.2) & Aplicación Android \\
multer & 1.4.5 & Recepción de imágenes \\
nodemailer & 8.0 & Envío de correo SMTP \\
firebase-admin & 14.3 & Notificaciones push (FCM) \\
nginx & 1.28.3 (alpine) & Servidor del frontend; el del host hace de proxy TLS \\
Docker Engine / Compose & $\geq$ 28 / v2 & Contenedores \\
Restic + Backblaze B2 & Restic 0.18.0 (en contenedor) & Copia externa cifrada \\
Jest & 29.7 & Pruebas unitarias de backend \\
\cod{node:test} & integrado en Node & Pruebas unitarias del frontend \\
\bottomrule
\end{tabularx}
\end{center}

\begin{nota}
La imagen de SQL Server del servidor está fijada a una actualización acumulativa concreta y a
su \emph{digest} SHA-256 en \ruta{deploy/server/compose.yaml}. Así, una reinstalación usa
exactamente la misma versión. El entorno de laboratorio usa la etiqueta \cod{2022-latest}.
\end{nota}

\section{Estructura del repositorio}
\begin{lstlisting}
servicedesk-microservices/
|-- api-gateway/        Gateway Express
|-- auth-service/       Servicio de autenticación (prisma/, src/, tests/)
|-- catalog-service/    Catálogo, ubicaciones y base de conocimiento
|-- ticket-service/     Tickets, adjuntos, reportes y notificaciones
|-- frontend/           React + Vite + Capacitor (android/)
|-- deploy/
|   |-- server/         compose.yaml y mssql-init/ del servidor nuevo
|   |-- nginx/          Plantillas de nginx del host (HTTP inicial y TLS)
|   |-- systemd/        Servicios y timers de backup y alertas
|   |-- backup/         Plantilla backup.env (credenciales de B2, fuera del .env)
|   |-- backup-lab/     Laboratorio MinIO para probar las copias externas
|   `-- certbot/, docker/, logrotate/, ssh/   Ajustes del host
|-- scripts/            Instalación, backups, operación y pruebas
|-- docs/               Documentación técnica
|-- archive/            Material retirado (PostgreSQL, frontend antiguo); no se usa
|-- compose.yaml        Entorno de laboratorio local (con .env.mssql)
|-- docker-compose.prod-mssql.yml   Compose del VPS anterior (solo referencia)
`-- .env.example        Plantilla de configuración (sin valores reales)
\end{lstlisting}
La carpeta \ruta{release/}, donde \cod{build-server-package.sh} deja los paquetes, no se versiona.

% =====================================================================
\chapter{Modelo de datos}\label{cap:datos}

Cada servicio tiene su base de datos y su usuario de SQL Server, con permisos solo
sobre esa base. Las relaciones entre servicios se hacen por identificador, por ejemplo
\cod{userId} o \cod{serviceId}, \emph{sin} claves foráneas entre bases. El esquema se
define en \ruta{<servicio>/prisma/schema.prisma}; existe también \ruta{schema.mssql.prisma}, idéntico,
que el \cod{Dockerfile.mssql} copia sobre el anterior al compilar la imagen.

\section{auth\_db}
\begin{center}
\begin{tabularx}{\linewidth}{@{}P{2.8cm}L@{}}
\toprule
\cabecera \textbf{Tabla} & \textbf{Campos principales} \\
\midrule
\cod{users} & id, email, name, role, microsoftId, passwordHash, mustChangePassword, block, createdAt, updatedAt \\
\cod{sessions} & id, userId, tokenHash (SHA-256 del JWT), authProvider, expiresAt, revokedAt, createdAt \\
\cod{auth\_logs} & id, userId, userEmail, userName, endpoint, authProvider, accessDate \\
\bottomrule
\end{tabularx}
\end{center}

\section{catalog\_db}
\begin{center}
\begin{tabularx}{\linewidth}{@{}P{3.6cm}L@{}}
\toprule
\cabecera \textbf{Tabla} & \textbf{Campos principales} \\
\midrule
\cod{services} & id, name, description, levelEntry, responsibleArea, isActive \\
\cod{routing\_rules} & id, serviceId, levelEntry, priorityHigh, priorityCritical, isCritical \\
\cod{locations} & id, name, isActive \\
\cod{knowledge\_articles} & id, title, problemDescription, solution, serviceId, createdByUserId, isActive \\
\cod{knowledge\_keywords} & id, articleId, keyword \\
\bottomrule
\end{tabularx}
\end{center}

\section{ticket\_db}
\begin{longtable}{@{}P{3.7cm}p{11.3cm}@{}}
\toprule
\cabecera \textbf{Tabla} & \textbf{Campos principales} \\
\midrule
\endhead
\cod{tickets} & id, number, userId, userName, userEmail, serviceId, serviceName, detail, location, status, priority, responsibleArea, assignmentStatus, assignedTechnicianId, assignedTechnicianName, acceptedAt, resolvedAt, resolvedById, resolvedByName, knowledgeArticleId, resolutionSummary, levelAssigned \\
\cod{ticket\_events} & id, ticketId, action, title, previousLevel/newLevel (heredados), previousStatus/newStatus, previousArea/newArea, previousTechnicianId/newTechnicianId, reason, workDone, description, visibility, performedBy, performedByName, createdAt \\
\cod{attachments} & id, ticketId, filePath, imageUrl, storageProvider, fileSize, mimeType, createdAt \\
\cod{acceptance\_terms} & version, title, text, sha256, createdAt \\
\cod{solution\_acceptances} & id, code, ticketId, decision, userId, userName, userEmail, userRole, decidedAt, termsVersion, termsSha256, solutionEventId, solutionSummary, technicianId, technicianName, reason \\
\cod{email\_outbox} & id, ticketId, to, subject, htmlBody, textBody, eventType, status, attempts, lastError, sentAt \\
\cod{device\_tokens} & id, userId, token, platform, deviceId, appVersion, isActive, lastSeenAt \\
\bottomrule
\end{longtable}

\begin{aviso}[Registros inmutables]
\cod{solution\_acceptances} y \cod{acceptance\_terms} están protegidas en la base de datos con
los disparadores \cod{trg\_solution\_acceptances\_immutable} y \cod{trg\_acceptance\_terms\_immutable}
(\cod{INSTEAD OF UPDATE, DELETE}), que ticket-service crea o actualiza en cada arranque. Ni la aplicación
ni un usuario con acceso a SQL pueden modificar o borrar las constancias ni los términos publicados.
Si el texto de una versión de los términos cambia, ticket-service deshabilita las aceptaciones
hasta que se publique una versión nueva.

\cod{ticket\_events} es un historial de \emph{solo inserción}: la aplicación nunca actualiza ni borra
eventos, y toda corrección se registra como un evento nuevo. Esta tabla \textbf{no} tiene disparador
de inmutabilidad en la base de datos.
\end{aviso}

\section{Valores de dominio}
\begin{center}
\begin{tabularx}{\linewidth}{@{}P{3.4cm}L@{}}
\toprule
\cabecera \textbf{Concepto} & \textbf{Valores} \\
\midrule
Estado del ticket & \cod{abierto}, \cod{en\_proceso}, \cod{escalado}, \cod{pendiente\_validacion}, \cod{resuelto}, \cod{cerrado} \\
Prioridad & \cod{baja}, \cod{media}, \cod{alta}, \cod{critica} \\
Área responsable & \cod{TECHNICIANS} (técnicos), \cod{DTIC} \\
Rol & \cod{admin}, \cod{technician}, \cod{dtic}, \cod{user} \\
Bloque & \cod{BLOQUE\_1}, \cod{BLOQUE\_2} (solo técnicos) \\
Estado de asignación & \cod{unassigned}, \cod{available}, \cod{accepted}, \cod{assigned\_by\_admin} \\
Acción de evento & \cod{CREATED}, \cod{CATEGORIZED}, \cod{AVAILABLE}, \cod{ACCEPTED}, \cod{ASSIGNED}, \cod{CONTRIBUTED}, \cod{STATUS\_CHANGED}, \cod{ESCALATED}, \cod{REASSIGNED}, \cod{COMMENTED}, \cod{RESOLVED}, \cod{USER\_ACCEPTED}, \cod{USER\_REJECTED\_RESOLUTION}, \cod{CLOSED} \\
Visibilidad de evento & \cod{public}, \cod{internal} \\
Decisión de validación & \cod{ACCEPTED}, \cod{REJECTED} \\
Estado de correo & \cod{pending}, \cod{sent}, \cod{failed} \\
Número de ticket & \cod{TK-AAAAMMDD-NNN}, consecutivo por día \\
Código de constancia & \cod{VD-AAAAMMDD-XXXXXXXX} (8 caracteres hexadecimales aleatorios) \\
\bottomrule
\end{tabularx}
\end{center}

\section{Aplicación del esquema y \emph{seed}}
El proyecto \textbf{no usa carpetas de migraciones}. Al arrancar, el \emph{entrypoint} de cada
microservicio (\ruta{<servicio>/docker/entrypoint.mssql.sh}) ejecuta \cod{prisma db push}, que
sincroniza la base con \ruta{schema.prisma} de forma idempotente:
\begin{itemize}
  \item En el servidor, \cod{PRISMA\_ACCEPT\_DATA\_LOSS=false}: si un cambio de esquema implicara borrar
        columnas o datos, el arranque \textbf{falla} en lugar de aplicarlo. Ese caso exige una migración manual
        planificada y un backup previo.
  \item En el laboratorio, la variable vale \cod{true} por defecto (\cod{--accept-data-loss}).
\end{itemize}

El \emph{seed} lo controla \cod{RUN\_SEED}. En el servidor vale \cod{false} y el instalador lo activa
\textbf{solo} en el primer arranque del modo limpio; en el modo de restauración nunca se siembra.
\begin{itemize}
  \item \textbf{catalog-service}: crea el catálogo de servicios, sus reglas, las ubicaciones y los
        artículos iniciales de conocimiento.
  \item \textbf{auth-service}: \textbf{no crea usuarios}. Las cuentas de prueba solo se crean con
        \cod{SEED\_DEMO\_USERS=true}, si existe el archivo local \ruta{prisma/seed.demo.ts} y \cod{NODE\_ENV}
        no es \cod{production}. El administrador inicial lo crea el instalador con \cod{manage-user.sh bootstrap-admin}.
  \item \textbf{ticket-service}: no tiene seed. Al arrancar publica la versión vigente de los términos
        de aceptación y crea los disparadores de inmutabilidad.
\end{itemize}

% =====================================================================
\chapter{Roles, áreas y bloques}

\begin{center}
\begin{tabularx}{\linewidth}{@{}P{2.2cm}P{2.9cm}L@{}}
\toprule
\cabecera \textbf{Rol} & \textbf{Áreas que atiende} & \textbf{Capacidades} \\
\midrule
\cod{user} & — & Crear tickets, ver los suyos y aceptar o rechazar la solución de sus tickets. \\
\cod{technician} & TECHNICIANS & Ver tickets disponibles de su área, aceptarlos, aportar, escalar y resolver. Pertenece a un bloque. \\
\cod{dtic} & DTIC & Lo mismo que el técnico, para los tickets escalados a la DTIC. \\
\cod{admin} & Ambas & Asignar tickets a técnicos, estadísticas, historial completo, reportes, catálogo, ubicaciones, base de conocimiento y bloques. No acepta tickets directamente: los asigna. \\
\bottomrule
\end{tabularx}
\end{center}

\begin{itemize}
  \item Los roles antiguos \cod{tech\_n1}/\cod{tech\_n2} se interpretan como \cod{technician}, y \cod{tech\_n3}/\cod{tech\_n4} como \cod{dtic}.
  \item Los \textbf{bloques} agrupan a los técnicos por zona física del campus y sirven para los reportes por bloque. Se asignan desde la pantalla \emph{Bloques} o con \cod{manage-user.sh set-block}.
  \item Un cambio de rol revoca las sesiones activas del usuario.
\end{itemize}

% =====================================================================
\chapter{Ciclo de vida del ticket}\label{cap:flujo}

\begin{figure}[h]
\centering
\begin{tikzpicture}[font=\small,
  est/.style={draw=utaRojo,rounded corners=8pt,minimum height=0.85cm,minimum width=2.6cm,align=center,fill=white},
  f/.style={-{Stealth[length=2.2mm]},line width=0.6pt},
  node distance=1.4cm]
  \node[est] (ab) {abierto};
  \node[est,right=of ab] (ep) {en\_proceso};
  \node[est,right=of ep] (pv) {pendiente\_\\validacion};
  \node[est,below=1.3cm of pv] (ce) {cerrado};
  \node[est,below=1.3cm of ep] (es) {escalado};
  \draw[f] (ab) -- node[above,font=\footnotesize]{aceptar} (ep);
  \draw[f] (ep) -- node[above,font=\footnotesize]{resolver} (pv);
  \draw[f] (pv) -- node[right,font=\footnotesize]{el creador acepta} (ce);
  \draw[f] (pv.north) to[bend right=40] node[above,font=\footnotesize]{el creador rechaza (con motivo)} (ep.north);
  \draw[f] ([xshift=-4mm]ep.south) -- node[left,font=\footnotesize]{escalar} ([xshift=-4mm]es.north);
  \draw[f] ([xshift=4mm]es.north) -- node[right,font=\footnotesize,align=left]{el área nueva\\acepta} ([xshift=4mm]ep.south);
\end{tikzpicture}
\caption{Estados principales del ticket.}
\end{figure}

El estado \cod{resuelto} se conserva por compatibilidad con tickets anteriores a la validación
digital: los tickets nuevos pasan de \cod{pendiente\_validacion} a \cod{cerrado}.

\begin{enumerate}
  \item \textbf{Creación}: el usuario elige el servicio, describe el problema, indica la ubicación
        (por ejemplo, un laboratorio) y adjunta hasta 5 imágenes JPG, PNG o WebP de 5\,MB como máximo.
        Las fotografías HEIC se rechazan con un mensaje que indica cómo cambiar el formato.
        El área responsable y la prioridad se toman de las reglas del catálogo.
  \item \textbf{Aceptación por el técnico}: el ticket aparece como \emph{disponible} para el personal del área.
        Un técnico lo acepta y queda \cod{en\_proceso}. El administrador puede asignarlo directamente
        a un técnico concreto.
  \item \textbf{Aportes}: el técnico registra el trabajo realizado sin cambiar el estado.
  \item \textbf{Escalamiento}: el técnico escala el ticket a la otra área con un motivo. El ticket pasa
        a \cod{escalado}, el técnico queda liberado y el ticket vuelve a la bandeja de disponibles del área
        nueva (aunque el área no cambie). Cuando alguien de esa área lo acepta, vuelve a \cod{en\_proceso}.
  \item \textbf{Resolución}: el técnico describe la solución y puede vincular un artículo de la
        base de conocimiento, cuyo nombre aparece en los reportes como «Solución utilizada».
        El ticket pasa a \cod{pendiente\_validacion}.
  \item \textbf{Validación digital}: \emph{solo} el usuario que creó el ticket puede decidir.
        \begin{itemize}
          \item \textbf{Acepta}: el ticket pasa a \cod{cerrado}.
          \item \textbf{Rechaza}: debe escribir un motivo. El ticket vuelve a \cod{en\_proceso} con el mismo técnico, que debe registrar una nueva solución. El ciclo se repite hasta que se acepte.
        \end{itemize}
        Cada decisión genera una constancia en \cod{solution\_acceptances} con un código único, la fecha,
        la versión de los términos aceptados y su huella SHA-256, y el resumen de la solución
        evaluada.
\end{enumerate}

Cada paso añade un evento a \cod{ticket\_events} dentro de la misma transacción que cambia el ticket.
El historial completo se consulta en la ventana de historial de cada ticket y en \emph{Historial}
del administrador, desde donde también se imprime.

% =====================================================================
\chapter{API REST}\label{cap:api}

Todas las rutas se publican bajo el dominio del sistema con el prefijo de la sección
\ref{sec:gateway}. Salvo que se indique lo contrario, requieren la cabecera
\cod{Authorization: Bearer <JWT>}. Las respuestas son JSON.

\begin{center}
\small
\begin{tabularx}{\linewidth}{@{}P{2.3cm}L@{}}
\toprule
\cabecera \textbf{Acceso} & \textbf{Significado} \\
\midrule
Público & No requiere sesión \\
Sesión & Cualquier usuario autenticado \\
Personal & \cod{admin}, \cod{technician} o \cod{dtic} \\
Técnico & \cod{technician} o \cod{dtic}, según el área del ticket \\
Admin & Solo \cod{admin} \\
Creador & Solo el usuario que creó el ticket \\
\bottomrule
\end{tabularx}
\end{center}

\section{Autenticación — \cod{/api/auth}}
\begin{longtable}{@{}P{1.7cm}P{5.3cm}P{1.9cm}p{5.6cm}@{}}
\toprule
\cabecera \textbf{Método} & \textbf{Ruta} & \textbf{Acceso} & \textbf{Descripción} \\
\midrule
\endhead
POST & \cod{/login} & Público & Inicio de sesión local. Limitado por cuenta y por IP \\
GET & \cod{/microsoft} & Público & Inicia el SSO de Microsoft Entra ID \\
GET & \cod{/microsoft/callback}, \cod{/callback} & Público & Retorno de Microsoft. Valida el \cod{state} y entrega un código de un solo uso \\
POST & \cod{/sso/exchange} & Público & Canjea el código de un solo uso por el JWT \\
GET & \cod{/config} & Público & Métodos de inicio de sesión disponibles \\
GET & \cod{/me} & Sesión & Datos del usuario actual \\
POST & \cod{/change-password} & Sesión & Cambio de contraseña. Revoca las demás sesiones \\
POST & \cod{/logout} & Sesión & Revoca la sesión actual \\
POST & \cod{/verify} & Sesión & Verifica el token (uso entre servicios) \\
POST & \cod{/users/lookup} & Personal & Datos básicos de hasta 200 usuarios por id \\
GET & \cod{/technicians} & Admin & Lista de técnicos (filtros \cod{role} y \cod{block}) \\
PATCH & \cod{/technicians/:id/block} & Admin & Asigna o retira el bloque de un técnico \\
POST & \cod{/dev-login}, \cod{/microsoft-simulate} & — & Solo desarrollo. En producción responden \textbf{404} \\
\bottomrule
\end{longtable}

\section{Catálogo — \cod{/api/catalog} y \cod{/api/knowledge}}
\begin{longtable}{@{}P{1.7cm}P{5.3cm}P{1.9cm}p{5.6cm}@{}}
\toprule
\cabecera \textbf{Método} & \textbf{Ruta} & \textbf{Acceso} & \textbf{Descripción} \\
\midrule
\endhead
GET & \cod{/api/catalog/services}, \cod{/services/:id} & Público & Servicios activos \\
POST\newline PATCH\newline DELETE & \cod{/api/catalog/services[/:id]} & Admin & Mantenimiento del catálogo \\
GET & \cod{/api/catalog/locations} & Público & Ubicaciones activas \\
POST\newline PUT\newline DELETE & \cod{/api/catalog/locations[/:id]} & Admin & Mantenimiento de ubicaciones \\
GET & \cod{/api/knowledge/search} & Personal & Búsqueda de artículos \\
GET & \cod{/api/knowledge/:id} & Personal & Detalle de un artículo \\
POST & \cod{/api/knowledge} & Personal & Crear un artículo \\
GET & \cod{/api/knowledge/}\allowbreak\cod{admin/list} & Admin & Lista completa, incluidos los inactivos \\
PUT\newline DELETE & \cod{/api/knowledge/:id} & Admin & Editar o desactivar un artículo \\
\bottomrule
\end{longtable}

\section{Tickets — \cod{/api/tickets}}
\begin{longtable}{@{}P{1.7cm}P{5.3cm}P{1.9cm}p{5.6cm}@{}}
\toprule
\cabecera \textbf{Método} & \textbf{Ruta} & \textbf{Acceso} & \textbf{Descripción} \\
\midrule
\endhead
POST & \cod{/} & Sesión & Crear un ticket (\cod{multipart/form-data} con imágenes) \\
GET & \cod{/my} & Sesión & Tickets creados por el usuario \\
GET & \cod{/:id} & Sesión & Detalle del ticket, si el usuario tiene acceso a él \\
GET & \cod{/history/:id} & Sesión & Historial de eventos \\
GET & \cod{/:id/attachments/}\allowbreak\cod{:attachmentId} & Sesión & Descarga autenticada de una imagen \\
GET & \cod{/available} & Personal & Tickets sin asignar de las áreas del usuario \\
POST & \cod{/:id/accept} & Técnico & Aceptar un ticket. El administrador debe usar la asignación \\
GET & \cod{/accepted}, \cod{/my-assigned} & Personal & Tickets aceptados o asignados \\
POST & \cod{/:id/contributions} & Técnico asignado & Registrar un aporte \\
PUT & \cod{/:id/escalate} & Técnico asignado & Escalar a la otra área con motivo \\
POST & \cod{/:id/resolve} & Técnico asignado & Registrar la solución \\
GET & \cod{/acceptance-terms} & Sesión & Texto vigente de los términos de aceptación \\
GET & \cod{/:id/acceptances} & Sesión & Constancias de aceptación o rechazo \\
POST & \cod{/:id/user-accept} & Creador & Aceptar la solución \\
POST & \cod{/:id/user-reject} & Creador & Rechazar la solución (motivo obligatorio) \\
\bottomrule
\end{longtable}

\section{Administración — \cod{/api/admin}}
\begin{longtable}{@{}P{1.7cm}P{5.3cm}P{1.9cm}p{5.6cm}@{}}
\toprule
\cabecera \textbf{Método} & \textbf{Ruta} & \textbf{Acceso} & \textbf{Descripción} \\
\midrule
\endhead
POST & \cod{/tickets/:id/assign} & Admin & Asignar un ticket a un técnico \\
GET & \cod{/stats/tickets} & Admin & Estadísticas generales \\
GET & \cod{/technicians}, \cod{/technicians/summary} & Admin & Técnicos y su carga \\
GET & \cod{/tickets/history} & Admin & Historial filtrable de tickets \\
GET & \cod{/reports/block} & Admin & Reporte por bloque y periodo \\
GET & \cod{/reports/technicians/}\allowbreak\cod{:technicianId} & Admin & Reporte por técnico y periodo \\
GET & \cod{/email-outbox} & Admin & Cola de correos \\
POST & \cod{/email-outbox/process} & Admin & Procesar la cola de correos manualmente \\
\bottomrule
\end{longtable}

\section{Notificaciones — \cod{/api/notifications}}
\begin{center}
\begin{tabularx}{\linewidth}{@{}P{1.7cm}P{5.3cm}P{1.9cm}L@{}}
\toprule
\cabecera \textbf{Método} & \textbf{Ruta} & \textbf{Acceso} & \textbf{Descripción} \\
\midrule
POST & \cod{/register}, \cod{/unregister} & Sesión & Registrar o retirar el token push del dispositivo \\
GET & \cod{/devices} & Sesión & Dispositivos registrados del usuario \\
\bottomrule
\end{tabularx}
\end{center}

\section{Rutas retiradas}
Las rutas del modelo anterior por niveles responden \textbf{410 Gone} con un mensaje que indica
qué usar en su lugar:
\cod{/api/routing/assign}, \cod{/api/routing/preview}, \cod{GET /api/tickets},
\cod{PUT /api/tickets/escalate}, \cod{GET /api/tickets/level} y
\cod{PATCH /api/tickets/:id/status}.

\section{Códigos de respuesta habituales}
\begin{center}
\begin{tabularx}{\linewidth}{@{}P{1.6cm}L@{}}
\toprule
\cabecera \textbf{Código} & \textbf{Situación} \\
\midrule
400 & Datos inválidos, archivo no permitido o código SSO inválido o caducado \\
401 & Sin token, token caducado o sesión revocada \\
403 & Rol sin permiso, o \cod{password\_change\_required} si el usuario debe cambiar su contraseña temporal \\
404 & Recurso inexistente, o ruta de desarrollo en producción \\
409 & Conflicto de estado (por ejemplo, aceptar un ticket ya tomado) \\
410 & Ruta retirada \\
429 & Demasiados intentos de inicio de sesión \\
\bottomrule
\end{tabularx}
\end{center}

% =====================================================================
\chapter{Frontend y aplicación móvil}

\section{Pantallas}
\begin{longtable}{@{}P{4.3cm}P{2.3cm}p{8.0cm}@{}}
\toprule
\cabecera \textbf{Ruta} & \textbf{Rol} & \textbf{Pantalla} \\
\midrule
\endhead
\cod{/login} & Público & Inicio de sesión (Microsoft o cuenta local) \\
\cod{/auth/callback}, \cod{/auth/success} & Público & Retorno del SSO y canje del código \\
\cod{/cambiar-contrasena} & Sesión & Cambio obligatorio de una contraseña temporal \\
\cod{/} & Sesión & Redirige al inicio que corresponde al rol \\
\cod{/tickets/nuevo} & Sesión & Crear un ticket \\
\cod{/mis-tickets} & Sesión & Tickets propios, con aceptación o rechazo de la solución \\
\cod{/tickets/:id} & Sesión & Enlace directo (por ejemplo, desde una notificación push). Redirige al listado del rol con \cod{?openTicket=<id>}, que abre el ticket \\
\cod{/panel} & Personal & Tickets disponibles \\
\cod{/mis-aceptados} & Personal & Tickets en atención \\
\cod{/conocimiento} & Personal & Base de conocimiento \\
\cod{/admin/stats} & Admin & Estadísticas \\
\cod{/admin/historial} & Admin & Historial de tickets e impresión \\
\cod{/admin/bloques} & Admin & Asignación de técnicos a bloques \\
\cod{/admin/reportes} & Admin & Reportes por bloque y por técnico \\
\cod{/admin/catalogo} & Admin & Catálogo de servicios \\
\cod{/admin/conocimiento} & Admin & Mantenimiento de artículos \\
\cod{/admin/ubicaciones} & Admin & Ubicaciones y laboratorios \\
\bottomrule
\end{longtable}

El detalle de un ticket (imágenes, historial y validación) se muestra en ventanas modales dentro de
los listados, no en una página propia. Cualquier otra ruta redirige a \cod{/}. Las imágenes se
comprimen en el navegador antes de subirlas y se muestran con una descarga autenticada.

\section{Reportes e impresión}
La pantalla \emph{Reportes} reúne en un solo formulario el reporte por \textbf{bloque} y el
reporte por \textbf{técnico}, con filtros de periodo y, para los bloques, de estado y servicio.
El documento conserva el formato institucional e incluye por cada ticket:
\begin{itemize}
  \item la solución utilizada y la ubicación o laboratorio;
  \item la validación digital del usuario (aceptado, rechazado o pendiente) con su código de constancia.
\end{itemize}

El documento se genera en \ruta{frontend/src/lib/institutionalReport.ts}. Para imprimir se usa un
\emph{iframe} oculto en la misma página, así que se abre un único diálogo de impresión sin ventanas
nuevas. Desde ese diálogo también se guarda en PDF. Las tablas tienen anchos de columna fijos y
evitan los cortes de fila entre páginas.

\section{Aplicación Android}
La aplicación se construye con Capacitor a partir del mismo frontend
(\ruta{frontend/android/}). Las notificaciones push usan Firebase Cloud Messaging. El archivo
\cod{google-services.json} y el almacén de claves de firma \textbf{no} se guardan en el repositorio.

\begin{bash}
cd frontend
npm ci && npm run build
npx cap sync android
# Abrir en Android Studio y generar el APK o AAB firmado
npx cap open android
\end{bash}

% =====================================================================
\chapter{Autenticación y seguridad}\label{cap:seguridad}

\section{Inicio de sesión}
\begin{itemize}
  \item \textbf{Microsoft Entra ID (SSO)}: es el método principal para la comunidad universitaria. Solo
        se admiten cuentas del dominio configurado en \cod{ALLOWED\_DOMAIN}. El servidor protege el retorno con un
        valor \cod{state} guardado en una cookie y no pasa el JWT en la URL. Entrega un código de un
        solo uso que caduca a los 60 segundos (se guarda solo su hash), y el frontend lo canjea en \cod{POST /sso/exchange}.
  \item \textbf{Cuenta local}: para el personal que lo necesite. Las contraseñas se guardan con
        \emph{bcrypt} (12 rondas). Los hashes SHA-256 heredados de versiones antiguas se aceptan una vez
        y se convierten a bcrypt tras el siguiente inicio de sesión correcto.
\end{itemize}

\section{Sesiones}
El JWT caduca según \cod{JWT\_EXPIRES\_IN} (8 horas por defecto). En la tabla \cod{sessions} se
guarda únicamente el hash SHA-256 de cada token. Así, una sesión se puede revocar en el cierre de sesión, en un cambio de
rol o de contraseña, y los registros caducados se eliminan pasados \cod{SESSION\_RETENTION\_DAYS} días (7 por defecto).
Como catalog-service y ticket-service consultan \cod{/auth/verify} en cada petición, la revocación
tiene efecto inmediato en todo el sistema.

\section{Política de contraseñas}
\begin{itemize}
  \item Siempre: mayúsculas, minúsculas, números y un carácter especial.
  \item Al menos \textbf{12 caracteres} al cambiar la contraseña o al crear una cuenta local con \cod{manage-user.sh}
        (la validación general admite 8, pero esos dos caminos exigen 12).
  \item No puede contener espacios al inicio o al final, ser igual al correo ni estar en la lista de contraseñas comunes.
  \item Las contraseñas temporales (administrador inicial y restablecimientos) obligan a cambiarlas en el primer acceso.
        Mientras tanto, la API responde \cod{403 password\_change\_required}.
\end{itemize}

\section{Protección contra fuerza bruta}
\begin{itemize}
  \item 5 fallos seguidos en una cuenta producen una espera progresiva de 1, 2, 4\ldots{} minutos, hasta un máximo de 15.
        Un acceso correcto reinicia el contador.
  \item 20 fallos desde una misma IP en 15 minutos bloquean esa IP durante 15 minutos.
  \item Antes de llegar a la aplicación, nginx del host limita \cod{/api/auth/login} a 10 peticiones por minuto por IP.
  \item El contador vive en la memoria de auth-service: se reinicia si el contenedor se reinicia.
  \item El mensaje de error es el mismo exista o no la cuenta. Los eventos se registran como
        \cod{[security]}, con el correo enmascarado y sin contraseñas.
\end{itemize}

\section{Otras medidas}
\begin{itemize}
  \item Las imágenes se sirven solo mediante una descarga autenticada, nunca como archivos estáticos públicos.
  \item Las rutas de desarrollo (\cod{dev-login}, \cod{microsoft-simulate}) responden 404 en producción.
        El \emph{seed} no crea cuentas. Las cuentas de prueba solo existen en un archivo local que Git no sigue.
  \item Cada servicio se conecta a SQL Server con un \emph{login} que solo tiene permisos sobre su base.
  \item En el host: UFW (22, 80, 443), SSH solo con clave, \cod{no-new-privileges} y rotación de logs en Docker.
        Ningún contenedor publica puertos en interfaces públicas.
  \item Los secretos solo existen en \ruta{deploy/server/.env} (permisos 600), que se genera en el servidor
        con \cod{--generate-secrets} y nunca se versiona.
\end{itemize}

% =====================================================================
\chapter{Notificaciones}
\begin{itemize}
  \item \textbf{Correo}: los eventos del ticket (creación, aceptación, asignación por el administrador, aporte,
        escalamiento, solución pendiente de validación y rechazo de la solución) encolan un mensaje en
        \cod{email\_outbox} (\cod{TICKET\_CREATED}, \cod{TICKET\_ACCEPTED}, \cod{TICKET\_ASSIGNED},
        \cod{TICKET\_CONTRIBUTION}, \cod{TICKET\_ESCALATED}, \cod{TICKET\_PENDING\_ACCEPTANCE},
        \cod{TICKET\_USER\_REJECTED}), que un proceso envía con reintentos. Con
        \cod{EMAIL\_SEND\_MODE=log}, los correos solo se escriben en el log, lo que sirve para pruebas. Con \cod{smtp} se
        envían mediante el servidor configurado en \cod{SMTP\_*}.
  \item \textbf{Push}: \cod{FCM\_MODE} admite \cod{off} (no envía nada), \cod{log} (solo registra) y
        \cod{firebase} (envía a los dispositivos registrados en \cod{device\_tokens}). Si se elige
        \cod{firebase} sin credenciales, el servicio pasa a \cod{log} y lo avisa en el log.
  \item \textbf{Alertas de operación}: los fallos de backup llaman a \cod{servicedesk-alert@} (\ruta{scripts/operations/alert.sh}), que
        escribe en el journal y en \ruta{/var/log/servicedesk-alerts.log}, y puede enviar un correo
        (\cod{BACKUP\_ALERT\_EMAIL}) o llamar a un \emph{webhook} (\cod{BACKUP\_ALERT\_WEBHOOK\_URL}).
        Se recomienda además un latido externo (\cod{BACKUP\_HEARTBEAT\_URL} en \ruta{backup.env}) que avise aunque el servidor entero caiga.
\end{itemize}

% =====================================================================
\chapter{Configuración}\label{cap:configuracion}

La configuración del servidor está en \ruta{deploy/server/.env}, creado a partir de
\ruta{.env.example}. Esta tabla lista las variables, \textbf{sin valores secretos}.

\begin{longtable}{@{}P{5.6cm}p{9.0cm}@{}}
\toprule
\cabecera \textbf{Variable} & \textbf{Uso} \\
\midrule
\endhead
\cod{PUBLIC\_DOMAIN} & Dominio público del sistema \\
\cod{PUBLIC\_SCHEME} & \cod{https} por defecto. Forma, con el dominio, la URL del frontend y de la API \\
\cod{GATEWAY\_BIND\_PORT}, \cod{FRONTEND\_BIND\_PORT} & Puertos locales (8080 y 3000), solo en 127.0.0.1 \\
\cod{SERVER\_PURPOSE} & \cod{pruebas} o \cod{produccion} (valor por defecto) \\
\cod{MSSQL\_PID} & Edición de SQL Server (obligatoria, sin valor por defecto). \cod{Developer} y \cod{Evaluation} quedan \textbf{bloqueadas} si \cod{SERVER\_PURPOSE=produccion} \\
\cod{MSSQL\_IMAGE} & Solo si se quiere sustituir la imagen fijada de SQL Server \\
\cod{MSSQL\_MEMORY\_LIMIT\_MB}, \cod{MSSQL\_CONTAINER\_MEMORY}, \cod{MSSQL\_CPUS} & Recursos de SQL Server \\
\cod{MSSQL\_SA\_PASSWORD} & Contraseña de \cod{sa} (generada) \\
\cod{AUTH\_DB\_*}, \cod{CATALOG\_DB\_*}, \cod{TICKET\_DB\_*} & Nombre, usuario y contraseña de cada base (contraseñas generadas) \\
\cod{JWT\_SECRET}, \cod{JWT\_EXPIRES\_IN} & Firma y duración de los tokens \\
\cod{MICROSOFT\_CLIENT\_ID}, \cod{\_TENANT\_ID}, \cod{\_CLIENT\_SECRET}, \cod{\_REDIRECT\_URI} & Registro de la aplicación en Microsoft Entra ID \\
\cod{ALLOWED\_DOMAIN} & Dominio de correo admitido \\
\cod{ADMIN\_EMAILS} & Correos que reciben el rol \cod{admin} al iniciar sesión con Microsoft \\
\cod{INITIAL\_ADMIN\_EMAIL}, \cod{INITIAL\_ADMIN\_NAME} & Administrador inicial (solo instalación limpia). No hay contraseña en el \cod{.env} \\
\cod{INITIAL\_ADMIN\_SSO\_ONLY} & \cod{true}: el administrador inicial solo entra con Microsoft (exige el SSO configurado) \\
\cod{SESSION\_RETENTION\_DAYS} & Días que se conservan las sesiones caducadas (7) \\
\cod{AUTH\_DEV\_LOGIN}, \cod{AUTH\_MICROSOFT\_SIMULATE}, \cod{SEED\_DEMO\_USERS} & Solo desarrollo. El compose del servidor las fija en \cod{false} \\
\cod{RUN\_SEED}, \cod{SERVICEDESK\_VERSION} & Los gestiona el instalador: seed del primer arranque y etiqueta de las imágenes \\
\cod{EMAIL\_NOTIFICATIONS\_ENABLED}, \cod{EMAIL\_SEND\_MODE}, \cod{SMTP\_*} & Correo \\
\cod{STORAGE\_MODE}, \cod{CLOUDINARY\_*} & Almacenamiento de imágenes (\cod{local} por defecto) \\
\cod{FCM\_MODE}, \cod{FCM\_PROJECT\_ID}, \cod{FCM\_CLIENT\_EMAIL}, \cod{FCM\_PRIVATE\_KEY} & Notificaciones push \\
\cod{BACKUP\_DIR}, \cod{BACKUP\_LOCAL\_KEEP}, \cod{BACKUP\_SQL\_COMPRESSION}, \cod{BACKUP\_REQUIRE\_OFFSITE} & Carpeta, conjuntos locales que se conservan, compresión y obligatoriedad de la copia externa \\
\cod{BACKUP\_ALERT\_EMAIL}, \cod{BACKUP\_ALERT\_WEBHOOK\_URL} & Destinos de las alertas de backup \\
\cod{DRILL\_MODE} & Modo del simulacro mensual de restauración \\
\cod{MSSQL\_HOST\_PORT}, \cod{API\_GATEWAY\_HOST\_PORT}, \cod{FRONTEND\_HOST\_PORT}, \cod{VITE\_API\_URL\_LAB} & Solo laboratorio local (\cod{compose.yaml}); el servidor no las usa \\
\bottomrule
\end{longtable}

\begin{aviso}
Las credenciales de Backblaze B2 y la contraseña del repositorio Restic \textbf{no} van en este
archivo. Se guardan aparte, en \ruta{/etc/servicedesk/backup.env}, con permisos 600.
\end{aviso}

% =====================================================================
\chapter{Requisitos del servidor}

\begin{center}
\begin{tabularx}{\linewidth}{@{}P{3.2cm}P{2.0cm}P{2.0cm}P{2.3cm}L@{}}
\toprule
\cabecera \textbf{Perfil} & \textbf{vCPU} & \textbf{RAM} & \textbf{Disco} & \textbf{Notas} \\
\midrule
Mínimo & 2 & 4 GB & 40 GB & Solo pruebas \\
Recomendado & 4 & 8 GB & 80 GB & Uso normal \\
Producción & 4 & 8 GB & 100 GB & SQL Server Express o edición con licencia \\
\bottomrule
\end{tabularx}
\end{center}

\begin{itemize}
  \item Ubuntu Server 24.04 LTS, arquitectura x86\_64 (SQL Server no funciona en ARM).
  \item Docker Engine 28 o superior con Compose v2, instalado desde el repositorio oficial de Docker.
  \item Un dominio con registro DNS hacia el servidor y los puertos 80 y 443 accesibles desde Internet para Let's Encrypt.
  \item En reposo, el sistema usa unos 1,6\,GiB de RAM.
  \item \textbf{Licencia}: la edición Developer solo puede usarse en servidores de pruebas. En producción se usa
        Express (límite de 10\,GB por base) o una edición con licencia.
\end{itemize}

% =====================================================================
\chapter{Instalación en un servidor nuevo}\label{cap:instalacion}

El sistema se distribuye como un paquete \cod{servicedesk-uta-server-<VERSION>.tar.gz}, acompañado de su
suma SHA-256. Se genera desde el repositorio con:
\begin{bash}
scripts/build-server-package.sh <VERSION>
\end{bash}
El resultado queda en \ruta{release/}: el \cod{.tar.gz}, su \cod{.sha256}, un \cod{MANIFEST.md} y un
\cod{FILES.txt} con la lista de archivos. El paquete excluye pruebas, documentación interna, archivos
\cod{.env} reales, \cod{node\_modules}, Android, copias de seguridad, imágenes subidas y cuentas de
prueba. Antes de empaquetar, el script aborta si encuentra secretos, claves privadas,
\cod{google-services.json}, valores de los \cod{.env} locales o referencias a PostgreSQL.

\begin{aviso}[Paquetes generados en Windows]
Un árbol de trabajo de Windows puede convertir los finales de línea a CRLF, y los scripts \cod{.sh}
fallan entonces en Linux. Genera el paquete desde un clon con \cod{core.autocrlf=false} y comprueba
los scripts en un contenedor Ubuntu antes de distribuirlo.
\end{aviso}

\section{Modos de instalación}
\begin{center}
\begin{tabularx}{\linewidth}{@{}P{3.6cm}L@{}}
\toprule
\cabecera \textbf{Modo} & \textbf{Cuándo usarlo} \\
\midrule
\cod{--mode clean} & Servidor nuevo sin datos. Crea las bases, ejecuta el seed del catálogo y crea el administrador inicial. \\
\cod{--mode restore --from DIR} & Servidor nuevo a partir de un backup de otro servidor. No ejecuta ningún seed y conserva usuarios, tickets e imágenes. \\
\bottomrule
\end{tabularx}
\end{center}

\section{Pasos}
\begin{enumerate}[leftmargin=*]
  \item \textbf{Verificar el paquete} en el servidor:
\begin{bash}
cd /tmp && sha256sum -c servicedesk-uta-server-<VERSION>.tar.gz.sha256
\end{bash}
  \item \textbf{Preparar Ubuntu} (actualizar el sistema y crear un usuario administrativo con clave SSH) e
        \textbf{instalar Docker Engine y Compose v2} desde el repositorio oficial.
  \item \textbf{Extraer} el paquete:
\begin{bash}
sudo mkdir -p /opt/servicedesk
sudo tar -xzf /tmp/servicedesk-uta-server-<VERSION>.tar.gz \
     -C /opt/servicedesk --strip-components=1
\end{bash}
  \item \textbf{Configurar} \ruta{deploy/server/.env}. Deja vacías las contraseñas y \cod{JWT\_SECRET}, porque el instalador las genera:
\begin{bash}
cd /opt/servicedesk
sudo cp .env.example deploy/server/.env && sudo chmod 600 deploy/server/.env
sudo nano deploy/server/.env
sudo scripts/install-server.sh --mode clean --generate-secrets --dry-run
\end{bash}
        El ensayo (\cod{--dry-run}) muestra cada variable como CONFIGURADA, VACÍA o AUSENTE, nunca su valor.
  \item \textbf{Instalar}:
\begin{bash}
# Instalación limpia
sudo scripts/install-server.sh --mode clean --generate-secrets
# o bien, restauración desde un backup
sudo scripts/install-server.sh --mode restore --from /root/restore/<ID> --generate-secrets
\end{bash}
  \item \textbf{DNS y HTTPS}: apuntar el dominio al servidor y obtener el certificado con certbot. El instalador
        deja configurado el \emph{hook} que recarga nginx en cada renovación.
  \item \textbf{Administrador inicial} (solo en el modo limpio): leer la contraseña temporal una vez y borrar el archivo:
\begin{bash}
sudo cat deploy/server/.initial-admin-secret
sudo shred -u deploy/server/.initial-admin-secret
\end{bash}
  \item \textbf{Comprobar}: iniciar sesión, crear un ticket con imagen y ejecutar \cod{server-health.sh}.
\end{enumerate}

\section{Qué hace el instalador}
\begin{enumerate}
  \item \textbf{Preflight}: sistema operativo, arquitectura, RAM, CPU, disco, Docker y Compose.
  \item \textbf{Configuración}: genera los secretos vacíos, valida las variables y pide confirmar la licencia de SQL Server.
  \item \textbf{Host}: instala \ruta{/etc/docker/daemon.json} y configura UFW. Endurece SSH \emph{solo} si ya existe una clave autorizada, y lo valida antes con \cod{sshd -t}. Instala también logrotate.
  \item \textbf{Compilación} de las 5 imágenes.
  \item \textbf{SQL Server}: lo arranca y \cod{mssql-init} crea las 3 bases con sus logins.
  \item \textbf{Servicios}: los arranca y espera \cod{6/6 servicios healthy}. En el modo limpio,
        el primer arranque lleva \cod{RUN\_SEED=true} y después los servicios se recrean sin seed.
  \item \textbf{Administrador inicial} (modo limpio): ejecuta \cod{manage-user.sh bootstrap-admin} con
        \cod{INITIAL\_ADMIN\_EMAIL} e \cod{INITIAL\_ADMIN\_NAME}.
  \item \textbf{Restauración} (modo restore): SHA-256, \cod{RESTORE VERIFYONLY}, \cod{DBCC CHECKDB}, imágenes, re-mapeo de usuarios de base y conteos.
  \item \textbf{Smoke test}: comprueba que las rutas de desarrollo respondan 404, las retiradas 410 y un código SSO inválido 400.
  \item \textbf{nginx} y \textbf{systemd}: configuración del proxy y timers de backup.
\end{enumerate}

\section{Opciones del instalador}
\begin{longtable}{@{}P{4.6cm}p{10.0cm}@{}}
\toprule
\cabecera \textbf{Opción} & \textbf{Efecto} \\
\midrule
\endhead
\cod{--mode clean|restore} & Modo de instalación \\
\cod{--from DIR} & Carpeta del backup (modo restore) \\
\cod{--env FILE}, \cod{--project NAME} & Archivo de configuración y nombre del proyecto de Compose \\
\cod{--generate-secrets} & Genera las contraseñas y el \cod{JWT\_SECRET} vacíos \\
\cod{--dry-run} & Valida sin cambiar nada \\
\cod{--yes} & No pide confirmaciones \\
\cod{--replace} & Solo en restore: permite sobrescribir bases existentes \\
\cod{--nginx-tls} & Instala la configuración de nginx con TLS (con el certificado ya emitido) \\
\cod{--skip-nginx}, \cod{--skip-systemd} & Omite nginx o los timers \\
\cod{--skip-firewall}, \cod{--skip-ssh-hardening}, \cod{--skip-docker-daemon} & Omite los ajustes del host \\
\cod{--ssh-port N} & Puerto SSH que UFW debe permitir (22 por defecto) \\
\cod{--test-mode} & Solo para la simulación local: omite las comprobaciones de root, sistema y arquitectura, y los ajustes de nginx, systemd, UFW, SSH y Docker. \textbf{No usar en un servidor real} \\
\bottomrule
\end{longtable}

El instalador no imprime secretos, no borra volúmenes, no ejecuta el seed sobre una base
restaurada y no modifica otras instalaciones. Nunca se debe ejecutar \cod{docker compose down -v},
porque borra los volúmenes con las bases y las imágenes.

% =====================================================================
\chapter{Operación}\label{cap:operacion}

En este capítulo, la abreviatura \cod{dc} equivale a:
\begin{bash}
docker compose -p servicedesk -f deploy/server/compose.yaml --env-file deploy/server/.env
\end{bash}

\section{Estado del sistema}
\begin{bash}
sudo scripts/operations/server-health.sh          # informe completo
sudo scripts/operations/server-health.sh --brief  # solo avisos y fallos
scripts/operations/smoke-server.sh                # prueba funcional rápida
curl -s http://127.0.0.1:8080/health/ready
dc ps
dc logs --tail 100 ticket-service
systemctl list-timers 'servicedesk-*'
\end{bash}

\cod{server-health.sh} revisa el host, los contenedores, SQL Server, la aplicación, la red, los backups,
los timers y los logs. Termina con código 0 si no hay fallos y con 1 si hay alguno. No muestra secretos.

\section{Gestión de usuarios}
\begin{longtable}{@{}P{6.6cm}p{8.0cm}@{}}
\toprule
\cabecera \textbf{Comando} (\ruta{scripts/operations/manage-user.sh}) & \textbf{Efecto} \\
\midrule
\endhead
\cod{bootstrap-admin <correo> ["<nombre>"] [--sso-only]} & Crea el primer administrador, solo si no existe ninguno, con una contraseña temporal guardada en \ruta{deploy/server/.initial-admin-secret} (permisos 600). Con \cod{--sso-only} no tiene contraseña local \\
\cod{create <correo> <rol> "<nombre>" [--password]} & Crea una cuenta que inicia sesión con Microsoft. Con \cod{--password} pide una contraseña local sin mostrarla (o la lee de \cod{NEW\_USER\_PASSWORD}) \\
\cod{set-role <correo> <rol>} & Cambia el rol y revoca las sesiones \\
\cod{set-block <correo> BLOQUE\_1|BLOQUE\_2|ninguno} & Asigna o retira el bloque de un técnico (también acepta \cod{NONE}) \\
\cod{reset-password <correo>} & Genera una contraseña temporal con cambio obligatorio y revoca sus sesiones \\
\cod{list-staff} & Lista el personal (rol y nombre, sin correos) \\
\bottomrule
\end{longtable}

\section{Actualizar la aplicación}
\begin{enumerate}
  \item Hacer un backup: \cod{sudo systemctl start servicedesk-backup.service}.
  \item Verificar y extraer el paquete nuevo sobre \ruta{/opt/servicedesk}, conservando \ruta{deploy/server/.env}.
  \item Recompilar y arrancar sin seed: \cod{RUN\_SEED=false dc up -d --build}. Las migraciones se aplican al arrancar.
  \item Ejecutar \cod{smoke-server.sh} y \cod{server-health.sh}.
  \item Si algo falla, volver al paquete anterior y, si la migración cambió datos, restaurar el backup del paso 1.
\end{enumerate}

% =====================================================================
\chapter{Copias de seguridad y recuperación}\label{cap:backups}

\section{Estrategia}
\begin{itemize}
  \item \textbf{Local}: \cod{backup-mssql.sh} genera un backup completo de las tres bases con \cod{BACKUP DATABASE}, junto con las imágenes y un \cod{SHA256SUMS}, en \ruta{/var/backups/servicedesk/sets/<ID>}.
  \item \textbf{Externa}: Restic cifra y envía cada conjunto a Backblaze B2. La clave del repositorio queda bajo la custodia de las personas designadas, y sin ella los backups externos no pueden recuperarse.
  \item \textbf{Verificación}: revisión periódica del repositorio y simulacro mensual de restauración en contenedores aislados.
\end{itemize}

\section{Tareas programadas (systemd)}
\begin{center}
\begin{tabularx}{\linewidth}{@{}P{5.6cm}P{3.6cm}L@{}}
\toprule
\cabecera \textbf{Timer} & \textbf{Frecuencia} & \textbf{Tarea} \\
\midrule
\cod{servicedesk-backup} & Diario, 03:15 & Backup local y envío a B2 \\
\cod{servicedesk-backup-check} & Cada 6 h (00:30, 06:30, 12:30, 18:30) & Comprueba la antigüedad y el estado del último backup \\
\cod{servicedesk-backup-retention} & Domingo, 04:30 & Aplica la política de retención \\
\cod{servicedesk-backup-verify} & Domingo, 05:30 & \cod{restic check} del repositorio \\
\cod{servicedesk-restore-drill} & Día 1 de cada mes, 06:00 & Simulacro de restauración \\
\bottomrule
\end{tabularx}
\end{center}

\section{Códigos de salida de \cod{backup-mssql.sh}}
\begin{center}
\begin{tabularx}{\linewidth}{@{}P{1.4cm}L@{}}
\toprule
\cabecera \textbf{Código} & \textbf{Significado} \\
\midrule
0 & Backup local y copia externa correctos \\
1 & Falló el backup local \\
3 & Backup local correcto, pero no hay copia externa configurada \\
4 & Backup local correcto, pero falló la copia externa \\
\bottomrule
\end{tabularx}
\end{center}

\section{Restaurar en el mismo servidor}
\begin{aviso}[Operación destructiva]
Sobrescribe las bases actuales. Haz antes un backup de lo que existe.
\end{aviso}
\begin{bash}
systemctl start servicedesk-backup.service                     # 0. Backup previo
scripts/restore-mssql.sh --from /var/backups/servicedesk/sets/<ID> --dry-run
dc stop auth-service catalog-service ticket-service api-gateway frontend
scripts/restore-mssql.sh --from /var/backups/servicedesk/sets/<ID> \
     --replace --with-uploads                                    # pide escribir RESTAURAR
RUN_SEED=false dc up -d
scripts/operations/smoke-server.sh
\end{bash}
\cod{restore-mssql.sh} verifica el SHA-256 y \cod{RESTORE VERIFYONLY} y comprueba que la versión del backup sea compatible.
Después restaura, ejecuta \cod{DBCC CHECKDB}, re-mapea los usuarios de base a los logins del servidor y muestra los conteos.

\section{Recuperar en un servidor nuevo}
Si el servidor original ya no existe, se recupera el conjunto más reciente desde B2 y después se
instala en modo restauración:
\begin{bash}
scripts/backup/restic-restore.sh --list
scripts/backup/restic-restore.sh --snapshot latest --target /root/restore/<ID>
sudo scripts/install-server.sh --mode restore --from /root/restore/<ID> --generate-secrets
\end{bash}
\cod{restic-restore.sh} solo restaura archivos en una carpeta nueva o vacía y verifica el contenido
descifrado, el \cod{SHA256SUMS} del conjunto y cada imagen. La carpeta contiene secretos
(\ruta{config/server.env}) y queda con permisos 700. Las contraseñas de SQL Server y el
\cod{JWT\_SECRET} serán nuevos. Los usuarios conservan sus cuentas y roles, pero deben volver a iniciar sesión.

% =====================================================================
\chapter{Desarrollo y pruebas}\label{cap:pruebas}

\section{Entorno local}
El laboratorio local levanta SQL Server, los microservicios, el gateway y el frontend con Docker,
usando \ruta{compose.yaml} en la raíz del repositorio:
\begin{bash}
cp .env.example .env.mssql      # MSSQL_SA_PASSWORD, *_DB_PASSWORD y JWT_SECRET
docker compose --env-file .env.mssql up -d --build
docker compose --env-file .env.mssql ps      # 6 servicios healthy
docker compose --env-file .env.mssql down    # nunca "down -v"
\end{bash}
Los puertos se publican solo en \cod{127.0.0.1}: 14330 (SQL Server), 18080 (gateway) y 13000 (frontend).
A diferencia del servidor, el laboratorio usa SQL Server Developer (\cod{2022-latest}),
\cod{NODE\_ENV=development} en auth-service (permite \cod{dev-login} y el inicio de sesión de Microsoft simulado),
correo en modo \cod{log}, almacenamiento local y FCM apagado.
Las cuentas de prueba se cargan desde \ruta{auth-service/prisma/seed.demo.ts}, un archivo local que
Git no sigue y que se crea a partir de \ruta{seed.demo.example.ts}. Ese archivo nunca se incluye en el paquete del servidor.

\section{Pruebas automatizadas}
\begin{center}
\begin{tabularx}{\linewidth}{@{}P{4.6cm}P{2.0cm}L@{}}
\toprule
\cabecera \textbf{Conjunto} & \textbf{Pruebas} & \textbf{Ejecución} \\
\midrule
auth-service & 76 & \cod{npm test} en la carpeta del servicio (Jest) \\
catalog-service & 71 & \cod{npm test} (Jest) \\
ticket-service & 270 & \cod{npm test} (Jest) \\
api-gateway & 3 & \cod{npm test} (Jest) \\
frontend & 31 & \cod{npm test} (\cod{node --test}) \\
\midrule
\textbf{Total unitarias} & \textbf{451} & \\
\bottomrule
\end{tabularx}
\end{center}
Recuento comprobado al ejecutar las cinco suites el 9 de noviembre de 2026: 451 de 451 correctas.
Antes de integrar cambios también debe compilar sin errores \cod{npx tsc --noEmit}.

\section{Pruebas de sistema}
\begin{longtable}{@{}P{5.6cm}p{9.0cm}@{}}
\toprule
\cabecera \textbf{Script} (\ruta{scripts/testing/}) & \textbf{Qué valida} \\
\midrule
\endhead
\cod{new-server-simulation.sh} & Instala el paquete en un servidor simulado y valida la instalación limpia, el smoke, la persistencia, la restauración y la caída de SQL Server \\
\cod{smoke-e2e.mjs} & Flujo funcional completo por la API: crear, aceptar, escalar, resolver, rechazar y aceptar \\
\cod{lab-copy-migration-test.sh} & Migración con una copia de los datos del laboratorio y verificación en el navegador \\
\cod{backup-dr-simulation.sh}, \cod{dr-validate.mjs} & Backup, recuperación ante desastre e integridad de los datos restaurados \\
\cod{browser-ui-check.mjs} & Pantallas en Chrome: imágenes, aceptación, rechazo e impresión de reportes \\
\bottomrule
\end{longtable}

% =====================================================================
\chapter{Solución de problemas}\label{cap:problemas}

\begin{longtable}{@{}P{5.0cm}p{9.6cm}@{}}
\toprule
\cabecera \textbf{Síntoma} & \textbf{Causa y solución} \\
\midrule
\endhead
El instalador rechaza \cod{MSSQL\_PID=Developer} & \cod{SERVER\_PURPOSE=produccion}. Usa \cod{Express} o la edición con licencia, o marca el servidor como \cod{pruebas} \\
\cod{mssql} nunca llega a \emph{healthy} & Poca RAM (mínimo 2\,GB para SQL Server) o permisos de los volúmenes. Revisa \cod{dc logs mssql} \\
Un servicio se reinicia en bucle & No conecta con su base. Comprueba que \cod{mssql-init} terminara bien y que las contraseñas del \cod{.env} no cambiaran tras la instalación \\
Error de CORS en el navegador & \cod{PUBLIC\_DOMAIN} no coincide con el dominio desde el que se accede \\
El SSO devuelve un error de \emph{redirect URI} & \cod{MICROSOFT\_REDIRECT\_URI} debe coincidir exactamente con la registrada en Entra ID (\cod{https://<dominio>/api/auth/microsoft/callback}) \\
Cuenta de Microsoft rechazada & El correo no pertenece a \cod{ALLOWED\_DOMAIN} \\
Respuesta 429 al iniciar sesión & Protección contra fuerza bruta. Espera el tiempo indicado o restablece la contraseña \\
Respuesta 403 \cod{password\_change\_required} & La cuenta tiene una contraseña temporal. Debe cambiarse desde la aplicación \\
Respuesta 410 & El cliente usa una ruta retirada. Actualiza el frontend o la aplicación \\
No se suben fotos del iPhone & Formato HEIC no admitido. Cambia la cámara a «Más compatible» o comparte la foto como JPG \\
Las imágenes no se ven tras restaurar & El backup no incluía \cod{uploads}. Restaura con \cod{--with-uploads} \\
Backup con código 3 & Falta configurar B2 en \ruta{/etc/servicedesk/backup.env} \\
Backup con código 4 & B2 no accesible o credenciales incorrectas. Usa \cod{restic-test-connection.sh} \\
El certificado caduca pronto & Revisa \cod{certbot renew --dry-run} y el puerto 80 \\
No llegan correos & \cod{EMAIL\_SEND\_MODE=log} o datos de \cod{SMTP\_*} incorrectos. Revisa \cod{/api/admin/email-outbox} \\
No llegan notificaciones push & \cod{FCM\_MODE=off}, credenciales de Firebase incorrectas o el dispositivo no está registrado \\
\bottomrule
\end{longtable}

La lista completa está en \ruta{docs/troubleshooting/PROBLEMAS_FRECUENTES.md}.

\section{Alcance de las pruebas realizadas}
Para interpretar correctamente este manual conviene distinguir qué se probó y dónde:
\begin{itemize}
  \item \textbf{Probado localmente}: pruebas unitarias, instalación del paquete en un servidor simulado
        con contenedores, flujo completo por la API, recuperación ante desastre con un repositorio
        Restic de laboratorio (MinIO) y comprobaciones en Chrome.
  \item \textbf{Aún no probado}: instalación en Ubuntu 24.04 real con systemd, nginx del host y certbot;
        SSO de Microsoft, envío SMTP y notificaciones FCM reales; copia externa en Backblaze B2;
        y el APK en un dispositivo físico.
\end{itemize}
Esos pasos deben validarse en el servidor de destino antes de ponerlo en producción.

% =====================================================================
\appendix
\chapter{Glosario}
\begin{description}[style=nextline,leftmargin=1.2cm]
  \item[API Gateway] Servicio que recibe todas las peticiones de la API y las reenvía al microservicio que corresponde.
  \item[Bloque] Agrupación de técnicos por zona del campus, usada en los reportes.
  \item[Constancia] Registro inmutable de la aceptación o el rechazo de una solución.
  \item[DBCC CHECKDB] Comprobación de integridad de una base de SQL Server.
  \item[Digest] Huella SHA-256 que identifica de forma exacta una imagen de contenedor.
  \item[Healthcheck] Comprobación periódica con la que Docker sabe si un contenedor está listo.
  \item[JWT] \emph{JSON Web Token}, credencial firmada que identifica una sesión.
  \item[Restic] Herramienta de copias de seguridad cifradas y deduplicadas.
  \item[Seed] Carga inicial de datos de referencia (catálogo, ubicaciones y conocimiento).
  \item[Smoke test] Prueba rápida de que las funciones esenciales responden.
  \item[SSO] \emph{Single Sign-On}, inicio de sesión con la cuenta institucional de Microsoft.
  \item[Validación digital] Aceptación o rechazo de la solución por parte del creador del ticket.
\end{description}

\end{document}
